% Part: counterfactuals
% Chapter: introduction
% Section: counterfactuals

\documentclass[../../../include/open-logic-section]{subfiles}

\begin{document}

\olfileid[pt]{cnt}{int}{cnt}

\olsection{Contrafactuais}

Uma forma muito comum e importante das construções inglesas ``se
\dots, então \dots'' é formada com o pretérito do subjuntivo do verbo
\emph{to be}: ``se fosse o caso que \dots, então seria o caso que
\dots''. Como o antecedente de tal condicional costuma ser falso, O.s.\@,
contrário aos fatos, esses condicionais são chamados de \emph{condicionais
  contrafactuais} (e, por usarem a forma subjuntiva de \emph{to
  be}, também de \emph{condicionais subjuntivos}). Eles se distinguem
dos condicionais \emph{indicativos}, que têm a forma ``se é o caso que
\dots, então é o caso que \dots''. Condicionais contrafactuais e
indicativos diferem quanto às condições de verdade. Considere o famoso
exemplo de Adams:
\begin{quote}
  Se Oswald não matou Kennedy, outra pessoa o matou.
  
  Se Oswald não tivesse matado Kennedy, outra pessoa o teria feito.
\end{quote}
O primeiro é indicativo; o segundo, contrafactual. O primeiro é
claramente verdadeiro: sabemos que o presidente John F. Kennedy foi morto por \emph{alguém}; se essa pessoa
não foi Lee Harvey Oswald (contrariamente ao Relatório Warren), então outra pessoa
matou Kennedy. O segundo diz algo diferente. Ele afirma
que, se Oswald não tivesse matado Kennedy, O.s.\@, se o tiroteio em Dallas tivesse
sido evitado ou fracassado, a história teria se desenrolado depois
de tal maneira que outro assassinato teria sido
bem-sucedido. Para que isso fosse verdadeiro, teria de ser o caso
que forças poderosas conspiraram para assegurar a morte de JFK (como acreditam muitos
teóricos da conspiração sobre JFK).

Ainda se debate se o condicional \emph{indicativo} é
representado corretamente pelo condicional material; em particular, se
os paradoxos do condicional material podem ser ``explicados'' de modo
compatível com o fato de ele fornecer as condições de verdade dos condicionais
indicativos do inglês. Em contraste, é incontroverso que
condicionais contrafactuais não podem ser simbolizados corretamente por
condicionais materiais. Isso é claro porque, embora os antecedentes
dos contrafactuais sejam geralmente falsos, nem todos os contrafactuais
com antecedentes falsos são verdadeiros---por exemplo, se acreditarmos no
Relatório Warren e não houve conspiração para assassinar JFK, então
o condicional contrafactual de Adams é um exemplo.

Condicionais contrafactuais desempenham um papel importante no
raciocínio causal: um exemplo central de seu uso é expressar
relações causais. P.ex., riscar um fósforo faz com que ele acenda, e
podemos expressar isso dizendo ``se este fósforo fosse riscado, ele
acenderia''. Condicionais materiais, e em geral os indicativos, não podem ser
usados para expressar isso: ``o fósforo é riscado $\lif$ o fósforo acende''
é verdadeiro se o fósforo nunca é riscado, independentemente do que aconteceria
se fosse. Pior ainda, ``o fósforo é riscado $\lif$ o fósforo se transforma
em um buquê de flores'' também é verdadeiro se ele nunca é riscado, mas
o fósforo certamente não se transformaria em um buquê de flores se
fosse riscado.

Ainda se debate qual é exatamente a lógica correta dos contrafactuais.
Stalnaker e Lewis apresentaram uma análise influente dos contrafactuais.
Segundo eles, um contrafactual ``se fosse o caso que~$S$, então seria o caso
que~$T$'' é verdadeiro se, e somente se, $T$ é verdadeiro na situação
contrafactual (``mundo possível'') mais próxima da maneira como o mundo real
é e na qual~$S$ é verdadeiro. Isso é chamado de análise
``ôntica'', pois faz referência a uma ontologia de
mundos possíveis. Outras análises usam probabilidades condicionais
ou teorias de revisão de crenças. Há uma proliferação de diferentes
lógicas propostas para os contrafactuais. Não há sequer uma única
lógica Lewis--Stalnaker dos contrafactuais: embora Stalnaker e
Lewis tenham proposto explicações semelhantes, com referência aos mundos possíveis
mais próximos, as suposições que fizeram resultam em diferentes inferências
válidas.

\end{document}
